\begin{table}[tbp]
\centering
\caption{问题二：共表达网络与稳定性选择}
\label{tab:t05_q2_network}
\footnotesize
\setlength{\tabcolsep}{3.5pt}
\begin{tabular}{lll}
\toprule
层面 & 指标 & 数值 \\
\midrule
网络 & 化合物节点数 & 223 \\
网络 & 共表达边数 & 1629 \\
网络 & 平均度 & 14.61 \\
网络 & 模块数 & 8 \\
网络 & 模块度 & 0.568 \\
筛选 & Lasso 稳定性选择 α & 0.06 \\
筛选 & Bootstrap 次数 B & 300 \\
筛选 & 分半次数 & 5 \\
稳定性 & 稳定率阈值 & 0.60 \\
稳定性 & 低稳定维度数 & 24/50 \\
稳定性 & 分半 Jaccard 中位数 & 1.0000 \\
图平滑 & 排序权重 γ & 0.20 \\
图平滑 & 分半 Jaccard（γ = 0–0.5 范围） & 1.000–1.000 \\
图平滑 & Top-3 Jaccard vs γ = 0（γ = 0.5） & 0.898 \\
\bottomrule
\end{tabular}
\end{table}
% 注（写入图注/表注，勿直接编译）：主表只收正向激动剂（Lasso positive、coef > 0、Spearman > 0）；负向掩蔽物另表。稳定性为 5 次随机分半、两半各自 Top-3 的 Jaccard；γ 事先固定为 0.20，未按测试结果调节。
